%% kapa: exam=1 round=36 year=2025 period=1 form=A subject=economics subject_ko=경제학원론 seq=20 number=60
%% source: Q-Net 감정평가사 기출문제 (cst00309) artlSeq=5231525 file=raw/1008/5231525/2025년 제36회 감정평가사 1차시험 1교시 문제지 원본.pdf
%% conversion: kordoc 4.24.51 -> tools/gen1.py (auto); review=proofread-sonnet; flags=none
\begin{question}{60}
X재의 수요곡선은 우하향하는 직선이고, 가로축 절편이 $a$ 이다. 공급곡선은 $\frac{a}{2}$
에서 수직인 직선이라고 할 때, 시장균형에서의 설명으로 옳은 것을 모두 고른 것은? (단, $a>0$ )
\begin{framedtext}
ㄱ. 수요의 가격탄력성은 절대치로 1이다. \par
ㄴ. 공급이 증가하면 기업의 총수입이 감소한다. \par
ㄷ. 공급이 증가하면 수요의 가격탄력성이 커진다. \par
ㄹ. 생산자잉여는 0이다.
\end{framedtext}
\begin{choices}
  \item ㄱ, ㄴ
  \item ㄱ, ㄷ
  \item ㄴ, ㄷ
  \item ㄴ, ㄹ
  \item ㄷ, ㄹ
\end{choices}
\end{question}
\answer{1}
